\chapter{Neue Drucksachen}
\quellseite{10}

Predigt nach dem großen Brand, gehalten am Sonntage Estomihi, den 14. Februar 1904, über 1 Petr. 5, 6. 7. vor der ev.-luth. St. Pauls-Gemeinde zu Baltimore, Md., und auf Verlangen dem Druck überlassen von Chr. Kühn, Pastor an genannter Gemeinde. 15 Seiten 7½ × 5. Zu beziehen vom Concordia Publishing House, St. Louis, Mo. Preis: 5 Cts.

Es war ohne Zweifel ganz richtig, daß unsere Pastoren am Sonntage nach dem schrecklichen Brande, der kürzlich Baltimore heimgesucht hat und der die Herzen aller Einwohner tief bewegte und erschütterte, diese schwere Heimsuchung Gottes zum Gegenstand einer besonderen Predigt machten und ihren Gemeinden dieses Ereigniß in das Licht des Wortes Gottes stellten und ihnen zeigten, was Gott durch dieses Unglück der ganzen Stadt, unserm ganzen Lande und besonders auch seinen Christen sagen und predigen will. Eine solche Predigt ist auf Wunsch der betreffenden Gemeinde dem Druck übergeben worden und liegt hier vor. Diese Predigt ist lehrreich und erbaulich nicht nur für unsere Gemeindeglieder in Baltimore, sondern auch für Christen aller Orten. Sie lernen daraus, wie sie sich zu verhalten haben, wenn Gott nach seinem weisen Rath sie einmal schwer heimsucht und züchtigt. Die Predigt hat 1 Petr. 5, 6. 7. zum Text, und auf Grund dieser Stelle zeigt der Verfasser: „Wie wir in dieser schweren Heimsuchung, die unsere Stadt getroffen hat, Gott die Ehre geben sollen. Das sollen wir thun, 1. indem wir bußfertig unter die gewaltige Hand Gottes uns demüthigen, und 2. indem wir nicht verzagen, sondern seiner väterlichen Fürsorge, die er in Christo uns verheißen hat, uns getrösten.“ G. M.

Verhandlungen der achten Versammlung der ev.-luth. Districts-Synode von Michigan. 56 Seiten 8½ × 6. Zu beziehen von Rev. C. Bast, Kawkawlin, Bay Co., Mich. Preis: 10 Cts.

Dieser Synodalbericht unserer Schwestersynode enthält außer den üblichen Jahresberichten und Geschäftsverhandlungen als Hauptgegenstand eine lehrreiche Abhandlung über das Verhalten der Zuhörer gegen ihre Prediger. An der Hand der Haustafel wird mit eingehender Behandlung der darin angeführten Schriftworte gezeigt, was die Zuhörer ihren Lehrern und Seelsorgern schuldig sind, nämlich sie im Tischreden zu versorgen, sie in Ehren zu halten, sie lieb zu haben und friedsam mit ihnen zu sein, ihnen zu gehorchen und sie nicht zu betrüben. Wenn einerseits die Pfarrerherren und Prediger, andererseits die Gemeindeglieder recht lernen und immer beherzigen, was ihnen in der Haustafel aus Gottes Wort gesagt wird, wird es gut in der Gemeinde stehen. L. F.

Der geplante Kirchenbund. Ein Wort der Warnung von Theodor Groß, Licentiaten der Theologie und Pastor zu Wetter in Hessen. Zwickau i. S. Druck und Verlag von Johannes Herrmann. 1903. 24 Seiten 9 × 6. Zu beziehen vom Concordia Publishing House, St. Louis, Mo. Preis 15 Cts.

Seit gerade zwei Jahren, als der deutsche Kaiser bei einer Festlichkeit die Vereinigung aller Evangelischen in Deutschland als ein Ziel seines Lebens bezeichnete, beschäftigt man sich drüben ernstlich mit der Sache, und letzten Sommer ist auf einer großen Kirchenconferenz die Sache auch so weit gediehen, daß die Vertreter fast sämtlicher deutsch-evangelischer Landeskirchen einen Plan entworfen und angenommen haben, der nur noch der endgültigen Annahme Seitens der einzelnen Kirchenregierungen bedarf. Gegen diesen Unionismus wendet sich der Verfasser der vorliegenden Schrift und weist ausführlich und lehrreich nach, warum sich ernste lutherische Christen an diesem geplanten Kirchenbund nicht betheiligen dürfen. Seine Schrift hat natürlich zunächst für deutschländische Christen Bedeutung, wird aber auch hier von allen denen mit Nutzen gelesen werden, die sich für die kirchlichen Zustände ihres alten Vaterlandes interessiren und über diese neueste unionistische Bewegung etwas Näheres wissen möchten. L. F.

Unentbehrlich für jedermann! Verlag des Schriftenvereins der sep. ev.-luth. Gemeinden in Sachsen. Zu beziehen vom Concordia Publishing House, St. Louis, Mo. 12 Seiten 7 × 5. Preis: 3 Cts.

Ein sehr guter Tractat. Während so viele Dinge als „unentbehrlich“ angepriesen werden, die doch oft sehr leicht entbehrt werden können, wird hier etwas genannt, was wirklich für jedermann unentbehrlich ist. Das ist die wahre Buße, die vor der Hölle bewahrt und in den Himmel führt. In schlichter, aber ansprechender Weise wird aus der Schrift gezeigt, was die wahre Buße ist, nämlich Reue über die Sünde und Glaube an Christum, und wie man zur wahren Buße kommt, nämlich allein durch Gottes Wort in Gesetz und Evangelium. Der Tractat eignet sich gut zur Verbreitung in der Stadtmission. L. F.

Die Religion der Odd Fellows oder der Sonderbaren Brüder. Von M. J. Von der Au, Route No. 1, Ida Grove, Iowa. 15 Seiten 7 × 5. Preis: 10 Cts., das Dutzend 75 Cts. portofrei.

Eine kleine Schrift, die gute Dienste leisten kann, wo man es mit der genannten weitverbreiteten Loge und mit solchen, die in ihre Netze verstrickt sind, zu thun hat. Aus dem eigenen Handbuch der Odd Fellows wird gezeigt, was sie von der Vaterschaft Gottes und der Bruderschaft der Menschen, von Sünde und Gerechtigkeit, von Himmel und Hölle lehren und wie ihre Lehre mit der heiligen Schrift stimmt. Das Ergebnis der Untersuchung ist, daß diese Loge, wenn sie sich auch mit noch so schönen Namen schmückt, doch in jedem der genannten Punkte der heiligen Schrift widerspricht, und daß darum ein Christ, der Gottes Wort seines Fußes Leuchte und ein Licht auf seinem Wege sein lassen und seinem Herrn und Heiland treu bleiben will, mit diesem Orden nichts zu schaffen haben darf. L. F.

Geistliche Lieder für Männerchöre. 5. Heft. Concordia Publishing House, St. Louis, Mo. 12 Seiten 11 × 8. Preis: 20 Cts.; Dutzendpreis: \$1.50 und Porto.

Das vorliegende Heft dieser vorzüglichen Sammlung bietet folgende Ostergesänge: Ich bin die Auferstehung und das Leben (Mauffl); Der Herr ist auferstanden (Mendelssohn); Man singet mit Freuden vom Sieg (Gläser); Bleib mir nah (Berner). Alle Gesänge sind gut; besonders aber ist der Mendelssohnsche Chor schwungvoll, triumphirend, so recht den christlichen Osterlauben zum Ausdruck bringend, zugleich mit einer feinen Orgelbegleitung versehen. Sehr ansprechend ist auch die letzte Nummer, die sich auch bei einem Begräbniß gut verwenden läßt. L. F.

Verzeichniß der Scheine und Oster-Katalog des Concordia Publishing House, St. Louis, Mo. 32 Seiten 7½ × 6.

Dieser neue Katalog unseres Verlagshauses wird allen denen ein guter Rathgeber sein, die nach passenden Geschenken für die Confirmations- und Osterzeit fragen, und wird auf Verlangen an irgend eine Adresse frei gesandt. L. F.
